\subsection{Technical Results}\label{sec:sobolev-technical-results}

\subsubsection{Density Theorems}\label{subsec:sobolev-density-theorems}

In this subsection, we shall show that $C^{\infty}$-functions with bounded support are dense in $H^k(\Omega)$. No assumptions on boundary regularity are needed. The same proofs work for $W^{k,p}(\Omega)$ if $p<\infty$. We first show that functions with bounded support are dense; of course, this is only of interest if $\Omega$ is unbounded.

\begin{Lemma}
\label{lem:sobolev-bounded-support-density}
Functions of bounded support are dense in $H^k(\Omega)$.
\end{Lemma}

\begingroup
\postdisplaypenalty=10000
\renewcommand{\endProof}{\par\nobreak\hfill$\blacksquare$\par\vspace{0.1cm}}
\begin{Proof}
Let $\phi\in\mathcal{D}(\mathbb{R}^m)$ be a function which equals $1$ for $|\bm{x}|\leq1$, $0$ for $|\bm{x}|\geq2$, and which takes values in between otherwise. For any $u\in H^k(\Omega)$, consider the sequence $u_n(\bm{x})=u(\bm{x})\phi_n(\bm{x})$, where $\phi_n(\bm{x})=\phi(\bm{x}/n)$. We first show that $u_n\to u$ in $L^2(\Omega)$. Let $\varepsilon>0$ be given and let $f$ be a continuous function such that $\|u-f\|_2\leq\varepsilon$. We find
\begin{equation}\label{eq:sobolev-bounded-support-cutoff-convergence}
\|u\phi_n-u\|_2\leq\|(u-f)\phi_n\|_2+\|u-f\|_2+\|f-f\phi_n\|_2\leq2\varepsilon+\|f-f\phi_n\|_2.
\end{equation}
It is easy to see that $f\phi_n$ converges to $f$ as $n\to\infty$. Hence $u\phi_n$ converges to $u$ in $L^2$. The convergence of derivatives is obtained by using the product rule of differentiation and a straightforward bootstrap argument.
\end{Proof}
\endgroup

\begin{Theorem}
\label{thm:sobolev-interior-smooth-density}
$C^{\infty}(\Omega)\cap H^k(\Omega)$ is dense in $H^k(\Omega)$.
\end{Theorem}

\begingroup
\postdisplaypenalty=10000
\renewcommand{\endProof}{\par\nobreak\hfill$\blacksquare$\par\vspace{0.1cm}}
\begin{Proof}
Let $\{\phi_i\}_{i\in\mathbb{N}}$ be a partition of unity on $\Omega$, subordinate to a locally finite covering $\{U_i\}_{i\in\mathbb{N}}$. Then for any $u\in H^k(\Omega)$, we have $u=\sum_{i\in\mathbb{N}}u\phi_i$ (in the sense of distributions, \emph{not} in the sense of convergence in $H^k(\Omega)$!). Moreover, $u\phi_i$ is in $H^k(\Omega)$, and we claim it is actually in $H_0^k(\Omega)$. Indeed, $u\phi_i$ can be extended by zero outside $\Omega$; this yields a function in $H^k(\mathbb{R}^m)$. By Corollary~\ref{cor:full-space-test-functions-sobolev-density}, there is a sequence $f_n\in\mathcal{D}(\mathbb{R}^m)$ such that $f_n\to u\phi_i$ in $H^k(\mathbb{R}^m)$. Let now $\psi$ be a test function which equals $1$ on the support of $\phi_i$ and $0$ outside $U_i$; then we have $f_n\psi\to u\phi_i\psi=u\phi_i$, and $f_n\psi\in\mathcal{D}(\Omega)$.

For every $i\in\mathbb{N}$, we can therefore choose $u_i^{(n)}\in\mathcal{D}(\Omega)$ such that
\begin{equation}\label{eq:sobolev-smooth-density-local-approximation-error}
\|u_i^{(n)}-u\phi_i\|_{k,2}\leq\dfrac{1}{n2^i}.
\end{equation}
It follows that
\begin{equation}\label{eq:sobolev-smooth-density-summed-approximation}
\sum_{i\in\mathbb{N}}u_i^{(n)}-u\phi_i=\sum_{i\in\mathbb{N}}u_i^{(n)}-u
\end{equation}
is in $H^k(\Omega)$, with norm $\leq1/n$. Hence $u^{(n)}:=\sum_{i\in\mathbb{N}}u_i^{(n)}$ converges to $u$ in $H^k(\Omega)$, and because of local finiteness $u^{(n)}$ is of class $C^{\infty}$.
\end{Proof}
\endgroup

Hence any function in $H^k(\Omega)$ can be approximated by functions whose derivatives exist in the classical sense. This is often useful in proofs. One can do everything for smooth functions first and then argue the general case by taking limits. However, Theorem~\ref{thm:sobolev-interior-smooth-density} is still unsatisfactory. In many cases one would like to have approximation by functions which are smooth up to the boundary, rather than just in the interior of $\Omega$. This cannot be done without some assumptions on the boundary of $\Omega$. The result we give now is not optimal, but will suffice for our purposes.

\begin{Lemma}
\label{lem:sobolev-extension-boundary-smooth-density}
Assume that $\Omega$ has the $k$-extension property. Then functions in $C^{\infty}(\overline\Omega)$ with bounded support are dense in $H^k(\Omega)$.
\end{Lemma}

The proof follows immediately from Corollary~\ref{cor:full-space-test-functions-sobolev-density}.
